\begin{table}[!h]
\caption{Effects on Blank Votes and Voter Turnout}
\renewcommand{\arraystretch}{1.25}
\setlength{\tabcolsep}{6pt}
\centering
\begin{tabular}[t]{lcccc}
\hline
 & \multicolumn{2}{c}{Share of blank votes} & \multicolumn{2}{c}{Voter turnout} \\
\cmidrule(l){2-3} \cmidrule(l){4-5}
  & (1) & (2) & (3) & (4)\\
\hline
Spanish share 1936-1978$\times$Post & -2.288*** & -3.847*** & -0.015 & 0.200\\
 & (0.553) & (0.689) & (0.242) & (0.327)\\

\addlinespace
Observations & 4680 & 4680 & 4680 & 4680\\
R$^2$ & 0.448 & 0.449 & 0.741 & 0.741\\
Municipality FE & Yes & Yes & Yes & Yes\\
Time FE & Yes & Yes & Yes & Yes\\
Election type FE & Yes & No & Yes & Yes\\
General elections only & No & Yes & No & No\\
Spanish share census & 1970 & 1980 & 1970 & 1980\\
\hline
\end{tabular}
\vspace{0.3em}
\captionsetup{font=footnotesize, justification=justified, singlelinecheck=false}
\caption*{\footnotesize Notes: The dependent variable in columns (1) and (2) is the share of blank votes over total voters. In columns (3) and (4), the dependent variable is voter turnout. Spanish share is defined as the share of Spanish-born immigrants who arrived between 1936 and 1978 over the total municipal population. Columns (1) and (3) use the 1970 census to construct this variable; columns (2) and (4) use the 1980 census. Standard errors clustered at the municipality level in parentheses (312 clusters). * $p<0.10$, ** $p<0.05$, *** $p<0.01$.}
\end{table}
